\documentclass[11pt]{article}
\usepackage[utf8]{inputenc}
\usepackage[spanish]{babel}
\usepackage{amsmath}
\usepackage{graphicx}
\usepackage{listings}
\usepackage{xcolor}
\usepackage{hyperref}
\usepackage{geometry}
\usepackage{booktabs}
\geometry{margin=2.5cm}

\lstset{
    basicstyle=\ttfamily\small,
    breaklines=true,
    frame=single,
    numbers=left,
    numberstyle=\tiny,
    keywordstyle=\color{blue},
    commentstyle=\color{gray},
    stringstyle=\color{teal}
}

\title{Sistema de Animación por Computador de un Océano \\ \large Proyecto Final 1 --- Computación Gráfica}
\author{Rivas Huanca Diego Raul \\ Alvarez Astete Jheeremy Manuel \\ Universidad Nacional de San Agustín de Arequipa \\ Escuela Profesional de Ciencia de la Computación}
\date{\today}

\begin{document}
\maketitle

\begin{abstract}
Este informe describe el diseño e implementación de un sistema de animación
por computador que simula la superficie de un océano en movimiento mediante
OpenGL. Se detalla la generación de la malla, el cálculo de alturas basado en
la superposición de olas sinusoidales, el cómputo analítico de normales, el
texturizado, la iluminación, la cámara, los elementos adicionales de la
escena (barco, isla, faro, skybox), la simulación de flotación real del
barco sobre la superficie deformada y un evento dinámico de \emph{ola
errática} que perturba la escena periódicamente.
\end{abstract}

\tableofcontents
\newpage

\section{Introducción}
Este trabajo presenta la implementación de un sistema de animación por
computador que simula la superficie de un océano en movimiento. El proyecto
fue desarrollado en C++20 usando OpenGL Core Profile (versión 3.3), GLFW
para la gestión de ventana/entrada, GLAD para la carga de funciones OpenGL,
ImGui para controles en tiempo de ejecución y CMake como sistema de build.

El objetivo principal fue desplazar el cálculo de alturas y normales de la
superficie al GPU (vertex shader), de modo que sea viable simular la
superposición de 20 olas sinusoidales en tiempo real sobre una malla
razonablemente densa (del orden de $150\times150$ vértices), y que además
otros elementos de la escena (en particular un barco) reaccionen a la forma
real del agua mediante la misma fórmula, evaluada en CPU en puntos
puntuales.

\section{Generación de la malla}
La superficie del océano se representa como una malla regular de triángulos
sobre el plano $xz$. La malla se construye en \texttt{Oceano::construirMalla(...)}
(\texttt{src/Oceano.cpp}) y se sube \textbf{una sola vez} a la GPU. Cada
vértice guarda únicamente la posición base $(x,0,z)$ y las coordenadas UV;
la componente $y$ (altura) y la normal se recalculan en el vertex shader en
cada frame, por lo que no es necesario volver a subir el buffer de vértices
mientras el océano anima.

Detalles de implementación:
\begin{itemize}
    \item \textbf{Resolución}: configurable al construir el objeto
    (\texttt{filas} $\times$ \texttt{columnas}), 150$\times$150 por defecto.
    \item \textbf{Separación entre puntos}: parámetro \texttt{separacion}
    pasado al \texttt{Builder}.
    \item \textbf{Índices}: se generan dos triángulos por celda de la
    grilla y se suben como \emph{Element Buffer Object} (EBO), dibujados
    con \texttt{glDrawElements}.
\end{itemize}

\begin{lstlisting}[language=C,caption={Generación de vértices e índices en Oceano::construirMalla}]
for (int fila = 0; fila < filas; ++fila) {
    for (int col = 0; col < columnas; ++col) {
        PuntoMalla p{};
        p.x = col * separacion - anchoTotal * 0.5f;
        p.y = 0.0f; // altura base; se recalcula en el shader
        p.z = fila * separacion - largoTotal * 0.5f;
        p.s = float(col) / (columnas - 1);
        p.t = float(fila) / (filas - 1);
        vertices.push_back(p);
    }
}
// dos triangulos por celda -> indices -> VBO + EBO + VAO
\end{lstlisting}

\subsection{Estructura de datos}
La estructura \texttt{PuntoMalla} (\texttt{src/PuntoMalla.h}) almacena
únicamente la posición base $(x,0,z)$ y las coordenadas UV $(s,t)$. No
guarda normal ni altura porque ambas se calculan analíticamente en el
vertex shader cada frame — esto reduce el ancho de banda CPU$\to$GPU y
evita tener que recalcular y resubir 22\,500 vértices por frame en CPU.

\section{Modelo matemático de las olas}
La altura de la superficie se modela como la superposición de $n$ olas
sinusoidales:
\begin{equation}
    h(x, z, t) = \sum_{i=1}^{n} A_i \cos\big(k_i (x\cos d_i + z\sin d_i) - 2\pi f_i t + p_i\big)
\end{equation}
con el número de onda derivado de la relación de dispersión de ondas de
gravedad en aguas profundas:
\begin{equation}
    k_i = \frac{4\pi^2 f_i^2}{g},\quad g = 9.81\ \mathrm{m/s^2}
\end{equation}

Parámetros de cada ola:
\begin{itemize}
    \item $A_i$ --- amplitud (metros)
    \item $d_i$ --- dirección de propagación (radianes)
    \item $f_i$ --- frecuencia (Hz)
    \item $p_i$ --- fase inicial
\end{itemize}

Los parámetros $A_i$, $d_i$, $f_i$ provienen de \texttt{data/spectrum.txt}
(3 columnas: amplitud, dirección en radianes, frecuencia), del cual se
cargan \textbf{20 olas}. El archivo no incluye la fase, por lo que se genera
aleatoriamente con semilla fija (\texttt{OlasDesdeArchivo}, semilla$=42$)
para que la animación sea reproducible entre ejecuciones.

\subsection{Cálculo en GPU vs. CPU}
Para sostener 20 olas sobre una malla de $\sim$150$\times$150 vértices a
60~fps, el cálculo de altura y normales se realiza en el vertex shader
(\texttt{shaders/oceano.vert}). El CPU solo prepara arrays de parámetros
(amplitud, dirección, frecuencia, número de onda y fase) y los sube como
\emph{uniforms}; el shader itera sobre las olas activas y suma su
contribución para cada vértice, en paralelo.

\begin{lstlisting}[language=C,caption={Fragmento de shaders/oceano.vert: altura y derivadas para la normal}]
for (int i = 0; i < numOlas; ++i) {
    float dirX = cos(direccion[i]);
    float dirZ = sin(direccion[i]);
    float k = numeroOnda[i];
    float faseTotal = k*(x*dirX + z*dirZ) - DOS_PI*frecuencia[i]*tiempo + fase[i];
    altura += amplitud[i] * cos(faseTotal);
    float derivComun = -amplitud[i] * k * sin(faseTotal);
    dHdx += derivComun * dirX;
    dHdz += derivComun * dirZ;
}
normal = normalize(vec3(-dHdx, 1.0, -dHdz));
\end{lstlisting}

\section{Cálculo de normales}
La normal de la superficie parametrizada por $(x, h(x,z), z)$ se obtiene a
partir de sus dos vectores tangentes:
\begin{align*}
\mathbf{t}_x &= (1,\ \partial h/\partial x,\ 0) \\
\mathbf{t}_z &= (0,\ \partial h/\partial z,\ 1)
\end{align*}
La normal es proporcional a $\mathbf{t}_z \times \mathbf{t}_x = (-\partial
h/\partial x,\ 1,\ -\partial h/\partial z)$, de ahí que en el shader se
normalice \texttt{vec3(-dHdx, 1.0, -dHdz)}.

Esta aproximación analítica es más precisa y más barata que estimar
normales promediando las de las caras vecinas (técnica usada en labs
anteriores con mallas estáticas): aquí la altura cambia cada frame, así
que promediar caras habría requerido recalcular vecindades constantemente,
mientras que la derivada analítica sale "gratis" del mismo cómputo que ya
se hace para la altura.

\section{Texturizado}
La textura de agua \texttt{textures/ocean.tga} (1024$\times$1024, TGA sin
comprimir) se carga con \texttt{stb\_image} en
\texttt{Oceano::cargarTextura(...)}. Las coordenadas UV se generan al
construir la malla, y en el fragment shader se aplica un desplazamiento
dependiente del tiempo para simular flujo:

\begin{lstlisting}[language=C,caption={Animación de UV en shaders/oceano.frag}]
vec2 uvAnimado = uv + vec2(tiempo * 0.02, tiempo * 0.015);
colorBase = texture(texturaOceano, uvAnimado).rgb;
\end{lstlisting}

En tiempo de ejecución la textura se puede activar/desactivar con la tecla
\texttt{T}.

\section{Iluminación y materiales}
Se implementó un modelo de iluminación Phong clásico en
\texttt{shaders/oceano.frag}, con las tres componentes estándar:
\begin{itemize}
    \item \textbf{Ambiente}: luz base constante que simula luz indirecta.
    \item \textbf{Difusa}: proporcional a $\max(\mathbf{n}\cdot\mathbf{l}, 0)$,
    da volumen a la superficie de las olas.
    \item \textbf{Especular}: $\big(\max(\mathbf{v}\cdot\mathbf{r}, 0)\big)^{\text{brillo}}$,
    con exponente de brillo $=64$ para simular el destello característico
    del agua.
\end{itemize}
Colores por defecto: ambiente $(0.15, 0.20, 0.25)$, difuso base
$(0.05, 0.25, 0.35)$, especular blanco.

Además del sol/luz principal, el \textbf{faro} añade una segunda fuente de
luz puntual cálida y parpadeante (ver \S\ref{sec:faro}), disponible para
extender el modelo a múltiples luces.

\section{Cámara}
Se utilizó una cámara orbital (\texttt{src/Camara.cpp}) parametrizada en
coordenadas esféricas (radio, yaw, pitch) respecto a un punto objetivo fijo
en el origen. Controles:
\begin{itemize}
    \item Rotar: click izquierdo + arrastrar el mouse.
    \item Zoom: scroll del mouse (ajusta el radio, con límites mínimo/máximo).
\end{itemize}

\section{Elementos adicionales}
Además del océano, la escena incluye cuatro elementos que se integran
visualmente con posición, escala e iluminación coherentes con la superficie
animada.

\subsection{Barco y flotación realista}
El \texttt{Barco} (\texttt{src/Barco.cpp}) es un objeto compuesto (casco,
mástil, vela) construido con primitivas procedurales. Su posición vertical
y su orientación se recalculan cada frame consultando directamente la
fórmula del océano en su posición $(x,z)$, sin ningún motor de física:

\begin{itemize}
    \item \textbf{Altura}: \texttt{Oceano::alturaEn(x,z)} evalúa
    $h(x,z,t)$ sumando las 20 olas (más la ola errática si está activa,
    ver \S\ref{sec:erratica}) en CPU, para un único punto.
    \item \textbf{Inclinación (roll y pitch)}: se calcula la pendiente
    local exacta del agua, \texttt{Oceano::gradienteEn(x,z)}, que devuelve
    $\partial h/\partial x$ y $\partial h/\partial z$ mediante la derivada
    analítica de cada ola (\texttt{Ola::gradiente}, la misma matemática que
    usa el vertex shader para la normal, evaluada aquí en un solo punto en
    vez de por vértice). El ángulo de balanceo lateral (\emph{roll}) es
    $\arctan(\partial h/\partial x)$ y el de cabeceo proa-popa
    (\emph{pitch}) es $\arctan(\partial h/\partial z)$; ambos se aplican
    como rotaciones (\texttt{Mat4::rotarZ} y \texttt{Mat4::rotarX}
    respectivamente) sobre la matriz de modelo del barco.
\end{itemize}

\begin{lstlisting}[language=C,caption={Barco::actualizar --- flotacion via pendiente analitica}]
alturaActual_ = oceano.alturaEn(posicionBase_.x, posicionBase_.z);

float dHdx = 0.0f, dHdz = 0.0f;
oceano.gradienteEn(posicionBase_.x, posicionBase_.z, dHdx, dHdz);

roll_  = std::atan(dHdx);
pitch_ = std::atan(dHdz);
\end{lstlisting}

Como la inclinación depende directamente de la pendiente real del agua bajo
el casco, el barco reacciona de forma consistente a cualquier ola —
incluida la ola errática— sin necesitar lógica adicional específica para
"volcarlo": es la misma flotación normal, simplemente sometida a una
pendiente mucho más pronunciada.

\subsection{Isla y rocas}
La isla y las rocas (\texttt{src/Isla.cpp}) se generan proceduralmente a
partir de esferas escaladas de forma no uniforme, con posición, escala y
tono de gris aleatorios pero de semilla fija, de modo que la geometría es
reproducible entre ejecuciones (útil para comparar capturas y grabar el
video de entrega).

\subsection{Faro}\label{sec:faro}
El faro (\texttt{src/Faro.cpp}) se modela como una torre cilíndrica más una
cúpula cónica. Además de decorativo, expone su posición
(\texttt{Faro::posicionLuz()}) como una segunda fuente de luz puntual, con
un parpadeo suave modulado por $0.7 + 0.3\sin(1.5\,t)$.

\subsection{Skybox}
En lugar de un cubemap con seis texturas externas, se renderiza una cúpula
(media esfera invertida) con un gradiente de color entre horizonte y cenit,
resuelto en el fragment shader según la altura relativa del vértice. Se
dibuja con \texttt{glDepthFunc(GL\_LEQUAL)} y sin escritura al depth buffer,
para que quede siempre al fondo de la escena.

\section{Ola errática}\label{sec:erratica}
Como evento dinámico adicional, la escena dispara periódicamente una
\textbf{ola errática}: una ola transitoria de amplitud considerablemente
mayor a las del espectro base (\texttt{amplitudMax}$=3.0$, frente a
$0.2$--$0.5$ de las olas normales), que aparece durante 6 segundos con una
envolvente temporal en forma de campana:
\begin{equation}
    E(u) = \sin(\pi u), \quad u = \frac{t - t_{\text{inicio}}}{\text{duración}} \in [0,1]
\end{equation}
de modo que la amplitud crece suavemente desde 0, alcanza su pico a mitad
del evento y vuelve a 0 sin transiciones bruscas.

La ola errática se implementa en \texttt{Oceano} como una ola "virtual"
adicional (\texttt{Oceano::olaErraticaActual()}) que se suma tanto a las
uniforms del shader (afecta visualmente la malla) como al cómputo en CPU de
\texttt{alturaEn()} y \texttt{gradienteEn()} (afecta al barco). Se activa
desde \texttt{Escena::actualizar()} con un intervalo aleatorio entre 18 y 35
segundos, en dirección también aleatoria, simulando una ola fuera de lo
común que puede hacer tambalear —o incluso inclinar de forma dramática— al
barco que esté cerca.

\section{Patrones de diseño aplicados}
\begin{itemize}
    \item \textbf{Strategy}: \texttt{FuenteOlas} / \texttt{OlasDesdeArchivo}
    / \texttt{OlasManual} permiten cambiar el origen de las olas sin
    modificar \texttt{Oceano}.
    \item \textbf{Builder}: \texttt{Oceano::Builder} permite construir el
    océano con una API encadenada y parámetros opcionales (tamaño,
    separación, textura, shaders, olas).
    \item \textbf{Factory Method}: \texttt{ElementoEscenaFactory} centraliza
    la creación de cada tipo de elemento adicional (barco, isla, faro,
    skybox), ocultando a \texttt{Escena} los detalles de construcción de
    cada uno.
    \item \textbf{Composición sobre herencia}: \texttt{ObjetoRenderizable}
    encapsula la subida a GPU y el dibujo de una malla; \texttt{Barco},
    \texttt{Isla} y \texttt{Faro} lo usan por composición en vez de heredar
    de una clase con lógica de OpenGL. \texttt{Escena} agrupa
    \texttt{Oceano}, la cámara, las luces y los elementos adicionales.
\end{itemize}

\section{Compilación y ejecución}
Requisitos:
\begin{itemize}
    \item Toolchain MSYS2 UCRT64 con GLFW instalado
    (\texttt{pacman -S mingw-w64-ucrt-x86\_64-glfw}).
    \item Dependencias vendoreadas en \texttt{external/} (\texttt{glad},
    \texttt{imgui-1.92.8} con backend \texttt{opengl3}) y
    \texttt{shared/stb\_image.h}.
\end{itemize}

Desde CLion: abrir el proyecto con el toolchain MSYS2 UCRT64 configurado y
compilar el target \texttt{cg\_oceano}. El \texttt{CMakeLists.txt} copia
automáticamente \texttt{shaders/}, \texttt{textures/} y \texttt{data/}
junto al ejecutable en cada build.

\textbf{Nota de compatibilidad:} en hardware gráfico más antiguo, la
inicialización de ImGui con \texttt{\#version 330} para el backend OpenGL3
puede fallar al compilar su shader interno. En ese caso, cambiar a
\texttt{ImGui\_ImplOpenGL3\_Init("\#version 120")} en \texttt{main.cpp}
resuelve el problema (con warnings inofensivos de "version deprecated").

\section{Controles}
\begin{itemize}
    \item Rotar cámara: click izquierdo + arrastrar
    \item Zoom: scroll del mouse
    \item Alternar textura: tecla \texttt{T}
    \item Alternar iluminación: tecla \texttt{L}
    \item Aumentar/reducir velocidad de animación: teclado numérico
    \texttt{+} / \texttt{-} (\texttt{KP\_ADD} / \texttt{KP\_SUBTRACT}), o el
    slider del panel ImGui
    \item Salir: \texttt{Esc}
\end{itemize}

\section{Resultados y capturas}

Vista general:
\begin{figure}[h!]
  \centering
  \includegraphics[width=0.8\textwidth]{images/captura_oceano.png}
  \caption{Vista general del océano (textura ON, iluminación ON).}
\end{figure}

\newpage
Textura:
\begin{figure}[h!]
  \centering
  \includegraphics[width=0.48\textwidth]{images/captura_textura_off.png}
  \includegraphics[width=0.48\textwidth]{images/captura_textura_on.png}
  \caption{Comparación textura apagada/encendida (tecla T).}
\end{figure}

Iluminación:
\begin{figure}[h!]
  \centering
  \includegraphics[width=0.48\textwidth]{images/captura_iluminacion_off.png}
  \includegraphics[width=0.48\textwidth]{images/captura_iluminacion_on.png}
  \caption{Comparación iluminación apagada/encendida (tecla L).}
\end{figure}

Ola errática y flotación del barco:
\begin{figure}[h!]
  \centering
  \includegraphics[width=0.8\textwidth]{images/captura_ola_erratica.png}
  \caption{Momento del evento de ola errática, con el barco inclinado sobre la pendiente pronunciada.}
\end{figure}

\section{Conclusiones}
Se logró desplazar el cómputo crítico (altura y normales) al GPU, permitiendo
simular la superposición de 20 olas en tiempo real sobre una malla de tamaño
moderado, y reutilizar la misma fórmula matemática en CPU para que otros
elementos de la escena (el barco) reaccionen de forma físicamente
consistente a la superficie, incluyendo eventos dinámicos como la ola
errática. Aprendizajes clave: manejo de arrays de uniforms en GLSL, cómputo
analítico de derivadas parciales tanto para normales (GPU) como para
pendiente local (CPU), y los \emph{trade-offs} entre precisión y
rendimiento al decidir qué se calcula dónde.

Mejoras futuras:
\begin{itemize}
    \item Espuma en las crestas (por ejemplo con mapas de ruido y un umbral
    sobre la curvatura de la superficie).
    \item Reflejos/refracciones mediante \emph{render-to-texture} y
    muestreo de entorno.
    \item Nivel de detalle (LOD) de la malla según distancia a la cámara.
    \item Casco del barco con geometría más detallada (proa afilada) o
    modelo \texttt{.obj} externo, en vez de primitivas procedurales.
\end{itemize}

\section{Referencias}
Enunciado del proyecto (Computación Gráfica, UNSA-EPCC); documentación
oficial de OpenGL/GLSL; documentación de ImGui y GLFW; Fournier \& Reeves
(1986), \emph{A simple model of ocean waves}, como referencia académica del
modelo de suma de senoides utilizado.

\end{document}

